% GENERATED FILE. Edit canonical lessons, then run npm run generate:books.
% canonical-source-hash: fnv1a64:42f23c8fe773c271
% canonical-lessons: JA-C168-kimaru, JA-C168-kimeru, JA-C168-kigen, JA-C168-nobasu, JA-C168-nobiru

\chapter{Deadlines and Putting Things Off}
\label{ch:ja-deadlines-and-putting-things-off}
\hlchaptermodality{168}

\begin{chapteropening}
\textbf{By the end of this chapter:} \emph{I can talk about deadlines and plans that move.}

Complete the last lesson of chapter 168: I can talk about deadlines and plans that move.
\end{chapteropening}

\section[\texorpdfstring{\ja{きまる} (kimaru)}{きまる (kimaru)}]{\ja{きまる} (kimaru) — to be decided; to be settled}

\label{lesson:JA-C168-kimaru}



\begin{quote}
Before the new one: say the Japanese for to become different, then the Japanese for calling off.
\end{quote}

\subsection*{You'll want to know: \ja{きまる}}
\textbf{\ja{きまる}} — \emph{kimaru} — \textquotedblleft{}to be decided; to be settled\textquotedblright{}.

\begin{quote}
\textbf{\ja{らいしゅう} \ja{よていが} \ja{きまります。}} — \emph{raishū yotei ga kimarimasu} — \textquotedblleft{}The plan will be decided next week.\textquotedblright{}
\end{quote}

Polite \textbf{\ja{きまります}}, \emph{kimarimasu}; a godan verb.

\begin{cousinweb}
\textbf{\ja{きまる}} (\emph{kimaru}) is a native Japanese word (wago), inherited rather than borrowed from Chinese; its kanji was attached later as a kun'yomi, and its older origin is not settled (JMdict, KANJIDIC2).
\end{cousinweb}

\begin{grammarlens}[title={What you've built}]
One more word for everyday life at A2.
\end{grammarlens}

\subsection*{Your turn}
\begin{itemize}
\raggedright
  \item \emph{Say it:} \emph{kimaru}
  \item \emph{Say it:} \emph{kimaru}, once more
  \item \emph{Say it:} \emph{raishū yotei ga kimarimasu} — The plan will be decided next week.
  \item \emph{Recall:} say the Japanese for to become different, then the Japanese for calling off, then say \emph{kimaru} again
\end{itemize}

\begin{tcolorbox}[breakable,colback=teal!4,colframe=teal!35!black,title={Before you move on}]
What does \ja{きまる} mean? (\textquotedblleft{}To be decided; to be settled\textquotedblright{}.) Say it once more.
\end{tcolorbox}
\section[\texorpdfstring{\ja{きめる} (kimeru)}{きめる (kimeru)}]{\ja{きめる} (kimeru) — to decide}

\label{lesson:JA-C168-kimeru}



\begin{quote}
Before the new one: say the Japanese for calling off, then the Japanese for to be decided.
\end{quote}

\subsection*{You'll want to know: \ja{きめる}}
\textbf{\ja{きめる}} — \emph{kimeru} — \textquotedblleft{}to decide\textquotedblright{}.

\begin{quote}
\textbf{\ja{あした} \ja{ははと} \ja{よていを} \ja{きめます。}} — \emph{ashita haha to yotei o kimemasu} — \textquotedblleft{}Tomorrow I will decide the plan with my mother.\textquotedblright{}
\end{quote}

Polite \textbf{\ja{きめます}}, \emph{kimemasu}; an ichidan verb. Its intransitive partner is \textbf{\ja{きまる}}, \emph{kimaru}.

\begin{cousinweb}
\textbf{\ja{きめる}} (\emph{kimeru}) is a native Japanese word (wago), inherited rather than borrowed from Chinese, and built on the same root as \emph{kimaru}, \textquotedblleft{}to be decided\textquotedblright{}; its kanji was attached later as a kun'yomi (JMdict, KANJIDIC2).
\end{cousinweb}

\begin{grammarlens}[title={What you've built}]
One more word for everyday life at A2.
\end{grammarlens}

\subsection*{Your turn}
\begin{itemize}
\raggedright
  \item \emph{Say it:} \emph{kimeru}
  \item \emph{Say it:} \emph{kimeru}, once more
  \item \emph{Say it:} \emph{ashita haha to yotei o kimemasu} — Tomorrow I will decide the plan with my mother.
  \item \emph{Recall:} say the Japanese for calling off, then the Japanese for to be decided, then say \emph{kimeru} again
\end{itemize}

\begin{tcolorbox}[breakable,colback=teal!4,colframe=teal!35!black,title={Before you move on}]
What does \ja{きめる} mean? (\textquotedblleft{}To decide\textquotedblright{}.) Say it once more.
\end{tcolorbox}
\section[\texorpdfstring{\ja{きげん} (kigen)}{きげん (kigen)}]{\ja{きげん} (kigen) — a deadline; a time limit}

\label{lesson:JA-C168-kigen}



\begin{quote}
Before the new one: say the Japanese for to be decided, then the Japanese for to decide.
\end{quote}

\subsection*{You'll want to know: \ja{きげん}}
\textbf{\ja{きげん}} — \emph{kigen} — \textquotedblleft{}a deadline; a time limit\textquotedblright{}.

\begin{quote}
\textbf{\ja{しゅくだいの} \ja{きげんは} \ja{あしたです。}} — \emph{shukudai no kigen wa ashita desu} — \textquotedblleft{}The homework deadline is tomorrow.\textquotedblright{}
\end{quote}

\begin{cousinweb}
\textbf{\ja{きげん}} (\emph{kigen}) is Sino-Japanese (kango): two Chinese characters read with their on'yomi, readings borrowed from Chinese. They are \emph{ki}, \textquotedblleft{}period, time\textquotedblright{} (Middle Chinese \emph{gi}) and \emph{gen}, \textquotedblleft{}limit\textquotedblright{} (\emph{heanX}) (KANJIDIC2; Middle Chinese in Baxter's transcription).
\end{cousinweb}

\begin{grammarlens}[title={What you've built}]
One more word for everyday life at A2.
\end{grammarlens}

\subsection*{Your turn}
\begin{itemize}
\raggedright
  \item \emph{Say it:} \emph{kigen}
  \item \emph{Say it:} \emph{kigen}, once more
  \item \emph{Say it:} \emph{shukudai no kigen wa ashita desu} — The homework deadline is tomorrow.
  \item \emph{Recall:} say the Japanese for to be decided, then the Japanese for to decide, then say \emph{kigen} again
\end{itemize}

\begin{tcolorbox}[breakable,colback=teal!4,colframe=teal!35!black,title={Before you move on}]
What does \ja{きげん} mean? (\textquotedblleft{}A deadline; a time limit\textquotedblright{}.) Say it once more.
\end{tcolorbox}
\section[\texorpdfstring{\ja{のばす} (nobasu)}{のばす (nobasu)}]{\ja{のばす} (nobasu) — to put off; to lengthen}

\label{lesson:JA-C168-nobasu}



\begin{quote}
Before the new one: say the Japanese for to decide, then the Japanese for a deadline.
\end{quote}

\subsection*{You'll want to know: \ja{のばす}}
\textbf{\ja{のばす}} — \emph{nobasu} — \textquotedblleft{}to put off; to lengthen\textquotedblright{}.

\begin{quote}
\textbf{\ja{らいしゅうの} \ja{よていを} \ja{のばします。}} — \emph{raishū no yotei o nobashimasu} — \textquotedblleft{}I will put off next week's plans.\textquotedblright{}
\end{quote}

Polite \textbf{\ja{のばします}}, \emph{nobashimasu}; a godan verb.

\begin{cousinweb}
\textbf{\ja{のばす}} (\emph{nobasu}) is a native Japanese word (wago), inherited rather than borrowed from Chinese; its kanji was attached later as a kun'yomi, and its older origin is not settled (JMdict, KANJIDIC2).
\end{cousinweb}

\begin{grammarlens}[title={What you've built}]
One more word for everyday life at A2.
\end{grammarlens}

\subsection*{Your turn}
\begin{itemize}
\raggedright
  \item \emph{Say it:} \emph{nobasu}
  \item \emph{Say it:} \emph{nobasu}, once more
  \item \emph{Say it:} \emph{raishū no yotei o nobashimasu} — I will put off next week's plans.
  \item \emph{Recall:} say the Japanese for to decide, then the Japanese for a deadline, then say \emph{nobasu} again
\end{itemize}

\begin{tcolorbox}[breakable,colback=teal!4,colframe=teal!35!black,title={Before you move on}]
What does \ja{のばす} mean? (\textquotedblleft{}To put off; to lengthen\textquotedblright{}.) Say it once more.
\end{tcolorbox}
\section[\texorpdfstring{\ja{のびる} (nobiru)}{のびる (nobiru)}]{\ja{のびる} (nobiru) — to be put off; to grow longer}

\label{lesson:JA-C168-nobiru}



\begin{quote}
Before the new one: say the Japanese for a deadline, then the Japanese for to put off.
\end{quote}

\subsection*{You'll want to know: \ja{のびる}}
\textbf{\ja{のびる}} — \emph{nobiru} — \textquotedblleft{}to be put off; to grow longer\textquotedblright{}.

\begin{quote}
\textbf{\ja{まつりは} \ja{らいしゅうに} \ja{のびます。}} — \emph{matsuri wa raishū ni nobimasu} — \textquotedblleft{}The festival will be put off until next week.\textquotedblright{}
\end{quote}

Polite \textbf{\ja{のびます}}, \emph{nobimasu}; an ichidan verb. Its transitive partner is \textbf{\ja{のばす}}, \emph{nobasu}.

\begin{cousinweb}
\textbf{\ja{のびる}} (\emph{nobiru}) is a native Japanese word (wago), inherited rather than borrowed from Chinese; its kanji was attached later as a kun'yomi, and its older origin is not settled (JMdict, KANJIDIC2).
\end{cousinweb}

\begin{grammarlens}[title={What you've built}]
One more word for everyday life at A2.
\end{grammarlens}

\subsection*{Your turn}
\begin{itemize}
\raggedright
  \item \emph{Say it:} \emph{nobiru}
  \item \emph{Say it:} \emph{nobiru}, once more
  \item \emph{Say it:} \emph{matsuri wa raishū ni nobimasu} — The festival will be put off until next week.
  \item \emph{Recall:} say the Japanese for a deadline, then the Japanese for to put off, then say \emph{nobiru} again
\end{itemize}

\begin{tcolorbox}[breakable,colback=teal!4,colframe=teal!35!black,title={Before you move on}]
What does \ja{のびる} mean? (\textquotedblleft{}To be put off; to grow longer\textquotedblright{}.) Say it once more.
\end{tcolorbox}
